\chapter*{概要}
\addcontentsline{toc}{chapter}{概要}

小型株の実験 S1〜S4 は、同じ 1,000 社（元の標本）で行った。この実験 R1 では、元の標本に含まれない新しい 1,000 社で S1〜S4 を同じコード・同じ設定で 1 回ずつやり直し（再現）、あわせて、S4 の形の BLM が比べる相手（範囲の全社、浮動株時価総額で加重）を上回るかを、新しい主な検定（R-BLM）として事前登録して確かめた。

\begin{keybox}[要点]
\begin{itemize}
\item \textbf{4 つの主な検定のうち \V{nreplicated} つが再現した}。決算発表文の読み取りが発表後 2〜60 日のドリフトを予測すること（S1：\V{rSonec}bp/SD、$t$ = \V{rSonet}、並べ替えの $p$ = \V{rSoneperm}）、発表後の変動率を予測すること（S2：\V{rStwoc}\%/SD、$t$ = \V{rStwot}）、文章のリスクの点数で共分散を補正しても TE の見積もりは良くならないこと（S3）は、新しい標本でも同じだった。TE の校正（S4）は再現しなかった。
\item \textbf{R-BLM は、事前登録の基準では\V{Bverdict}された}（2017 年以降、比べる相手との差は年 \V{Bex}、95\% 区間 [\V{Blo}, \V{Bhi}]、片側 $p$ \V{Bp}）。ただし実現の TE は \V{Bte} で、約束の 3\% の約 3 倍だった。
\item \textbf{探索（結果を見た後の分析）で、その原因は株価データの誤りだとわかった}。\V{rtopday} に、破産から再出発した会社（CHRD）の古い株と新しい株の価格を yfinance がつないでいて、1 日で約 +25{,}700\% の、実際にはなかった値上がりが入っていた。この 1 日だけで、BLM と相手の差が \V{rtopactive} 開き、そのあと約 3 年、校正が TE の上限を下限の 1\% に張り付かせた。
\item \textbf{ありえない日々のリターン（+300\% 超、$-$80\% 未満）を除いて計算し直すと}、新しい標本でも校正は元の標本と同じように効いた（偏り $g$：校正なし \V{rmainBg} → 校正あり \V{rmainBcalg}。元の標本は \V{omainBg} → \V{omainBcalg}）。S4 の不成立は、この種のデータの誤りでほぼ説明できる。
\item 誤りを除いた後の R-BLM の BLM は、新しい標本で実現の TE \V{rmainBcalte}、相手との差は年 \V{rmainBcalex}（95\% 区間 [\V{rmainBcallo}, \V{rmainBcalhi}]、片側 $p$ \V{rmainBcalp}）、元の標本で差は年 \V{omainBcalex}（片側 $p$ \V{omainBcalp}）だった。事前登録の値（年 \V{Bex}）の多くは誤りから来ていたが、誤りを除いても、両方の標本で相手を上回った。これは結果を見た後の分析なので、誤りを除く規則を先に決めて、別の期間で確かめ直す必要がある。
\item 同じ種類のデータの誤りは元の標本にもあった。これからの小型株の計算では、株価の誤りを先に取り除く手順を、事前登録に入れる。
\end{itemize}
\end{keybox}
